\section{Introduction}
\paragraph{Installation}
Create a file on your computer named ``main.c'',
inside your are going to write the following:
\begin{minted}[bgcolor=mygreen!10]{c}
  #include <stdio.h>

  int main(void) {
    printf("Hello Seamen!\n");
    return 0;
  }
\end{minted}
now open a \textit{terminal/shell} of your choice
to \textit{compile} and \textit{run} the ``main.c''
file with the following:
\begin{minted}[bgcolor=mygreen!10]{bash}
  # I am on Linux and use 'gcc' depending
  # on your operating system you may need
  # to use something else.
  gcc -o main.c main
  ./main
\end{minted}
if the phrase ``Hello Seamen!'' appeared on your
terminal its good, you can move to the next paragraph.
If nothing happened make sure everything has been properly
written, if an error occurs \textit{google} ``How to install
C compiler on Windows/Mac/Linux'' or ask a trusted friend to
do so for you. Repeat this process until the
desired result has been reached.

\paragraph{Pre-requisites}
In order to move to the next chapter it is required
to understand the behaviour of the following program
without any external help:
\begin{minted}[bgcolor=mygreen!10]{c}
  // 'int' is an abbreviation for integer
  // '=' is an assignment, translate it with 'set to'
  int factorial(int n) {
    int x = 1;
    loop:
    if (n > 0) {
      x = x * n;
      n = n - 1;
      goto loop;
    }
    else {
      return x;
    }
  }
\end{minted}
Can you guess the \textit{return} value of $x$ if $n = 5$?
If you understood the code correctly it should be $x = 120$.
